% !TEX root = task-1-18-2.tex
% Editable task bank. Append new BankTask blocks at the end of the document.
\documentclass[12pt,a4paper]{article}
\usepackage{physics-textbook}
\usepackage{graphicx,needspace}
\usetikzlibrary{decorations.pathmorphing}
\newif\ifShowSolutions
\ifdefined\TutorEdition
  \ShowSolutionstrue
\else
  \ShowSolutionsfalse % Change to \ShowSolutionstrue to display solutions.
\fi
\newcounter{banktask}
\newsavebox{\BankTaskBox}
\newcommand{\BankTaskContent}[3]{%
  \noindent\textbf{\thebanktask.}\quad #2\par
  \ifShowSolutions
    \Needspace{50mm}% Keep the solution heading with its opening text or diagram.
    #3\par
    \ifdefined\TutorEdition
      \ifcsname TutorSource#1\endcsname
        \nopagebreak\smallskip{\small\csname TutorSource#1\endcsname\par}%
      \fi
    \fi
  \fi
}
\newcommand{\BankTask}[3]{%
  \refstepcounter{banktask}\par\medskip
  \sbox{\BankTaskBox}{\begin{minipage}{\linewidth}%
    \BankTaskContent{#1}{#2}{#3}\end{minipage}}%
  \ifdim\dimexpr\ht\BankTaskBox+\dp\BankTaskBox\relax>\dimexpr\textheight-20mm\relax
    \Needspace{35mm}\BankTaskContent{#1}{#2}{#3}%
  \else
    \Needspace{\dimexpr\ht\BankTaskBox+\dp\BankTaskBox\relax}%
    \noindent\usebox{\BankTaskBox}\par
  \fi
  \medskip
}
\begin{document}
\BookSubsection{1.18.2}{Звук}


% Source sdamgia-640: https://phys-ege.sdamgia.ru/problem?id=640
\BankTask{sdamgia-640}{
Какова частота звуковых колебаний в среде, если скорость звука в этой среде $v =500\,\text{м/с},$ а длина волны $\lambda =2\,\text{м}$?
}{
\smallskip\noindent\textbf{Решение}\par
Частота, длина и скорость распространения звуковых волн в среде связаны соотношением $\lambda \nu = v _{\,\text{зв}}.$ Отсюда находим частоту звуковых колебаний:

$\nu = \frac{v _{\,\text{зв}}}{\lambda} = \frac{500\,\text{м/с}}{2\,\text{м}} =250\,\text{Гц}.$

Ответ: $250$.
}

% Source sdamgia-39917: https://phys-ege.sdamgia.ru/problem?id=39917
\BankTask{sdamgia-39917}{
Чему равна скорость звуковой волны в воздухе, если её длина равна $1{,}65\,\text{м}$, а частота колебаний $\nu = 200\,\text{Гц}?$
}{
\smallskip\noindent\textbf{Решение}\par
Скорость распространения волны равна

$v =\lambda \nu=1{,}65 \cdot 200=330 \,\text{м/с} .$

Ответ: $330$.
}

% Source sdamgia-621: https://phys-ege.sdamgia.ru/problem?id=621
\BankTask{sdamgia-621}{
Диапазон голоса мужского баса занимает частоты от $\nu_1 = 80\,\text{Гц}$ до $\nu_2 = 400\,\text{Гц}.$ Каково отношение граничных длин звуковых волн $\frac{\lambda_1}{\lambda_2}$ этого диапазона?
}{
\smallskip\noindent\textbf{Решение}\par
Частота, длина и скорость распространения звуковых волн связаны соотношением $\lambda \nu =c.$ Отсюда находим отношение граничных длин волн:

$\frac{\lambda _1}{\lambda _2} = \frac{c/\nu _1}{c/\nu _2} = \frac{\nu _2}{\nu _1} = \frac{400\,\text{Гц}}{80\,\text{Гц}} =5.$

Ответ: $5$.
}

% Source sdamgia-626: https://phys-ege.sdamgia.ru/problem?id=626
\BankTask{sdamgia-626}{
Человеческое ухо воспринимает звуковые волны, длины которых лежат в интервале от $\lambda_1=16\,\text{мм}$ до $\lambda_2=20\,\text{м}.$ Каково отношение граничных частот звуковых волн $\frac{\nu_1}{\nu_2}$ этого интервала?
}{
\smallskip\noindent\textbf{Решение}\par
Частота, длина и скорость распространения звуковых волн связаны соотношением $\lambda \nu =c.$ Отсюда находим отношение граничных частот волн:

$\frac{\nu_1}{\nu_2} = \frac{c/\lambda_1}{c/\lambda_2} = \frac{\lambda_2}{\lambda_1} = \frac{20\,\text{м}}{0{,}016\,\text{м}} =1250.$

Ответ: $1250$.
}

% Source sdamgia-642: https://phys-ege.sdamgia.ru/problem?id=642
\BankTask{sdamgia-642}{
На расстоянии $400\,\text{м}$ от наблюдателя рабочие вбивают сваи с помощью копра. Каково время между видимым ударом молота о сваю и звуком удара, услышанным наблюдателем? Скорость звука в воздухе $330\,\text{м/с}.$ Округлите ответ с точностью до десятых.
}{
\smallskip\noindent\textbf{Решение}\par
Можно считать, что свет преодолевает расстояние в $400$ метров практически мгновенно. Звуку на это требуется:

$\frac{400\,\text{м}}{330\,\text{м/с}} \approx 1{,}2\,\text{с}.$

Ответ: $1{,}2$.
}

% Source sdamgia-3355: https://phys-ege.sdamgia.ru/problem?id=3355
\BankTask{sdamgia-3355}{
Звуковой сигнал, отразившись от препятствия, вернулся обратно к источнику через $5\,\text{с}$ после его испускания. Каково расстояние от источника до препятствия, если скорость звука в воздухе $340\,\text{м/с}$?
}{
\smallskip\noindent\textbf{Решение}\par
За $5\,\text{с}$ звуковой сигнал успел дважды пройти расстояние между источником и препятствием, а значит, это расстояние равно:

$\frac{v t}{2} = \frac{340 \cdot 5}{2} =850$ м.

Ответ: $850$.
}

% Source sdamgia-643: https://phys-ege.sdamgia.ru/problem?id=643
\BankTask{sdamgia-643}{
Для экспериментального определения скорости звука ученик встал на расстоянии $30\,\text{м}$ от стены и хлопнул в ладоши. В момент хлопка включился электронный секундомер, который выключился отраженным звуком. Время, отмеченное секундомером, равно $0{,}18\,\text{с}$. Какова скорость звука, определенная учеником? Округлите до целых.
}{
\smallskip\noindent\textbf{Решение}\par
Звук проделал путь до стены и обратно перед тем, как выключить секундомер. Следовательно, скорость звука, определенная учеником, равна:

$v = \frac{2S}{t} = \frac{2 \cdot 30\,\text{м}}{0{,}18\,\text{с}} \approx 333\,\text{м/с}.$

Ответ: $333$.
}

% Source sdamgia-6640: https://phys-ege.sdamgia.ru/problem?id=6640
\BankTask{sdamgia-6640}{
Гидроакустик, находящийся на корабле, переговаривается по рации с матросом, находящимся на лодке. Во время разговора матрос наносит удар гаечным ключом по корпусу своей лодки. Звук от этого удара гидроакустик сначала слышит через рацию, а через $10$ секунд  — через свою гидроакустическую аппаратуру. Считая, что второй звук распространяется в воде со скоростью $1500\,\text{м/с}$, найдите расстояние между кораблем и лодкой.
}{
\smallskip\noindent\textbf{Решение}\par
Считая, что сигнал по рации приходит почти мгновенно, найдем расстояние, которое проходит звук от лодки до корабля: $1500\,\text{м/с} \cdot 10\,\text{с} = 15 000\,\text{м} = 15\,\text{км}.$

Ответ: $15$.
}

% Source sdamgia-24945: https://phys-ege.sdamgia.ru/problem?id=24945
\BankTask{sdamgia-24945}{
Две лодки покоятся на поверхности озера на некотором расстоянии друг от друга. С первой лодки одновременно посылаются два коротких звуковых сигнала, один из которых идет в воздухе, а второй  — в воде. На второй лодке один из сигналов был принят через $0{,}8\,\text{с}$ после отправки, а другой  — на $2{,}7\,\text{с}$ позже первого сигнала. Найдите отношение скорости звука в воде к скорости звука в воздухе. Ответ округлите до десятых долей.
}{
\smallskip\noindent\textbf{Решение}\par
Оба сигнала прошли одинаковое расстояние: $s= v _1t,$ $s= v _2 ( t + \Delta t ) .$ Тогда получаем отношение скорости звука в воде к скорости звука в воздухе: 

$\frac{v _1}{v _2} = \frac{t + \Delta t}{t} = \frac{0{,}8 + 2{,}7}{0{,}8} \approx 4{,}4.$

Ответ: $4{,}4$.
}

% Source sdamgia-11659: https://phys-ege.sdamgia.ru/problem?id=11659
\BankTask{sdamgia-11659}{
На поверхности моря покоится катер. Непосредственно под ним на глубине $50\,\text{м}$ работает водолаз, который в некоторый момент ударяет молотком по металлической детали. Сидящий на катере гидроакустик слышит два звука от удара с интервалом времени между ними $1\,\text{с}$. Скорость звука в воде $1400\,\text{м/с}$. Чему равна глубина моря в этом месте?
}{
\smallskip\noindent\textbf{Решение}\par
Обозначим за $h$ глубину, на которой работает водолаз, $x$ — расстояние от водолаза до дна моря, $H$ — глубину моря в данной точке. Тогда время прихода гидроакустику первой звуковой волны равно $t_1= \frac{h}{v} ,$ второй звуковой волны  — $t_2= \frac{2x + h}{v} .$ Таким образом, вторая волна пришла позже первой на промежуток времени $\Delta t=t_2 - t_1= \frac{2x}{v} .$
Тогда расстояние от водолаза до дна равно $x= \frac{v \Delta t}{2} ,$ а глубина моря:

$H=h + x=h + \frac{v \Delta t}{2} =50 + \frac{1400 \cdot 1}{2} =750\,\text{м}.$

Ответ: $750$.
}


% Source A784FF: https://ege.fipi.ru/bank/questions.php?search=1&proj=BA1F39653304A5B041B656915DC36B38&qid=A784FF
\BankTask{A784FF}{
Человек услышал звук грома через $8\,\text{с}$ после вспышки молнии. Считая, что скорость звука в воздухе равна $340\,\text{м}/\text{с}$, определите, на каком расстоянии от человека ударила молния. Результат выразите в $\text{м}$.
}{
\smallskip\noindent\textbf{Решение от GPT}\par
Время распространения света пренебрежимо мало. Расстояние $s=vt=340\cdot8=2720\,\text{м}$.
}

% Source 3F1DFE: https://ege.fipi.ru/bank/questions.php?search=1&proj=BA1F39653304A5B041B656915DC36B38&qid=3F1DFE
\BankTask{3F1DFE}{
У входа в вертикальную шахту произведён выстрел. Через какое время после выстрела звук выстрела вернётся к стрелку, отразившись от дна шахты, если её глубина $170\,\text{м}$? Скорость звука в воздухе принять равной $340\,\text{м}/\text{с}$. Результат выразите в $\text{с}$.
}{
\smallskip\noindent\textbf{Решение от GPT}\par
Звук проходит путь $2h$, поэтому $t=2h/v=2\cdot170/340=1\,\text{с}$.
}

% Source 991714: https://ege.fipi.ru/bank/questions.php?search=1&proj=BA1F39653304A5B041B656915DC36B38&qid=991714
\BankTask{991714}{
Труба (бас) издаёт звуки в диапазоне от $\nu_{1}=60\,\text{Гц}$ до $\nu_{2}=6000\,\text{Гц}$.
Чему равно отношение граничных длин звуковых волн $\frac{\lambda_1}{\lambda_2}$ этого диапазона?
}{
\smallskip\noindent\textbf{Решение от GPT}\par
В одной среде $\lambda=v/\nu$, поэтому $\lambda_1/\lambda_2=\nu_2/\nu_1=6000/60=100$.
}

% Source E92D12: https://ege.fipi.ru/bank/questions.php?search=1&proj=BA1F39653304A5B041B656915DC36B38&qid=E92D12
\BankTask{E92D12}{
Каков период колебаний $T$ звуковых волн в среде, если скорость звука
в этой среде равна $1000\,\text{м}/\text{с}$, а длина волны $l=5\,\text{м}$? Результат выразите в $\text{с}$.
}{
\smallskip\noindent\textbf{Решение от GPT}\par
$v=\lambda/T$, откуда $T=\lambda/v=5/1000=0{,}005\,\text{с}$.
}

% Source 4e8F5A: https://ege.fipi.ru/bank/questions.php?search=1&proj=BA1F39653304A5B041B656915DC36B38&qid=4e8F5A
\BankTask{4e8F5A}{
Звуковая волна частотой $2\,\text{кГц}$ распространяется в среде со скоростью $340\,\text{м}/\text{с}$. Чему равна длина звуковой волны? Результат выразите в $\text{м}$.
}{
\smallskip\noindent\textbf{Решение от GPT}\par
$\nu=2000\,\text{Гц}$. Длина волны $\lambda=v/\nu=340/2000=0{,}17\,\text{м}$.
}

% Source 39CA85: https://ege.fipi.ru/bank/questions.php?search=1&proj=BA1F39653304A5B041B656915DC36B38&qid=39CA85
\BankTask{39CA85}{
Точно посередине между наблюдателем и скалой в дерево ударила молния.
Наблюдатель услышал первый раскат грома через $1\,\text{с}$ после вспышки молнии. Через какое время после вспышки молнии наблюдатель услышит отражённый от скалы звук? Результат выразите в $\text{с}$.
}{
\smallskip\noindent\textbf{Решение от GPT}\par
Пусть расстояние до молнии $d$. Первый звук проходит $d$ за $1\,\text{с}$. До скалы от молнии также $d$, затем от скалы до наблюдателя $2d$; отражённый звук проходит $3d$ и придёт через $3\,\text{с}$ после вспышки.
}

\end{document}
